\SourceLeaf{104}
ist jetzt unter dem Namen Methode der kleinsten Quadrate im allgemeinen Gebrauch, und ihre Begründung durch die Wahrscheinlichkeitsrechnung, so wie die Bestimmung der Genauigkeit der Resultate selbst, nebst andern damit zusammenhängenden Untersuchungen sind in der \textit{Theoria Motus Corporum Coelestium} ausführlich entwickelt.

Der Marquis \textsc{Delaplace}, welcher nachher diesen Gegenstand aus einem neuen Gesichtspunkte betrachtete, indem er nicht die wahrscheinlichsten Werthe der unbekannten Grössen suchte, sondern die zweckmässigste Combination der Beobachtungen, fand das merkwürdige Resultat, dass, wenn die Anzahl der Beobachtungen als unendlich gross betrachtet wird, die Methode der kleinsten Quadrate allemal und unabhängig von der Function, die die Wahrscheinlichkeit der Fehler ausdrückt, die zweckmässigste Combination sei.

Man sieht hieraus, dass beide Begründungsarten noch etwas zu wünschen übrig lassen. Die erstere ist ganz von der hypothetischen Form für die Wahrscheinlichkeit der Fehler abhängig, und sobald man diese verwirft, sind wirklich die durch die Methode der kleinsten Quadrate gefundenen Werthe der unbekannten Grössen nicht mehr die wahrscheinlichsten, eben so wenig wie die arithmetischen Mittel in dem vorhin angeführten einfachsten aller Fälle. Die zweite Begründungsart lässt uns ganz im Dunkeln, was bei einer mässigen Anzahl von Beobachtungen zu thun sei. Die Methode der kleinsten Quadrate hat dann nicht mehr den Rang eines von der Wahrscheinlichkeitsrechnung gebotenen Gesetzes, sondern empfiehlt sich nur durch die Einfachheit der damit verknüpften Operationen.

Der Verfasser, welcher in gegenwärtiger Abhandlung diese Untersuchung aufs neue vorgenommen hat, indem er von einem ähnlichen Gesichtspunkt ausging, wie \textsc{Delaplace}, aber den Begriff des mittlern zu befürchtenden Fehlers auf eine andere, und wie ihm scheint, schon an und für sich natürlichere Art, feststellt, hofft, dass die Freunde der Mathematik mit Vergnügen sehen werden, wie die Methode der kleinsten Quadrate in ihrer neuen hier gegebenen Begründung allgemein als die zweckmässigste Combination der Beobachtungen erscheint, nicht näherungsweise, sondern nach mathematischer Schärfe, die Function für die Wahrscheinlichkeit der Fehler sei, welche sie wolle, und die Anzahl der Beobachtungen möge gross oder klein sein.

Mit dem Hauptgegenstande ist eine Menge anderer merkwürdiger Unter\-\SourceLeaf{105}suchungen enge verbunden, deren Umfang aber den Verfasser nöthigte, die Entwickelung des grössten Theils derselben einer künftigen zweiten Vorlesung vorzubehalten. Von denjenigen, die schon in der gegenwärtigen ersten Abtheilung vorkommen, sei es uns erlaubt, hier nur ein Resultat anzuführen. Wenn die Function, welche die relative Wahrscheinlichkeit jedes einzelnen Fehlers ausdrückt, unbekannt ist, so bleibt natürlich auch die Bestimmung der Wahrscheinlichkeit, dass der Fehler zwischen gegebene Grenzen falle, unmöglich: dessenungeachtet muss, wenn nur allemal grössere Fehler geringere (wenigstens nicht grössere) Wahrscheinlichkeit haben als kleinere, die Wahrscheinlichkeit, dass der Fehler zwischen die Grenzen \(-x\) und \(+x\) falle, nothwendig grösser (wenigstens nicht kleiner) sein, als \(\frac{x}{m}\surd\frac13\), wenn \(x\) kleiner ist als \(m\surd\frac34\), und nicht kleiner als \(1-\frac{4mm}{9xx}\), wenn \(x\) grösser ist als \(m\surd\frac34\), wobei \(m\) den bei den Beobachtungen zu befürchtenden mittlern Fehler bedeutet. Für \(x=m\surd\frac34\) fallen wie man sieht beide Ausdrücke zusammen.

\par\bigskip
\begin{center}
\rule{.88\textwidth}{.3pt}\par
Göttingische gelehrte Anzeigen.\quad 1823 Februar 24.\par
\rule{.88\textwidth}{.3pt}
\end{center}
\pdfbookmark[1]{1823 Februar 24}{anzeige-1823-02-24}

Eine am 2. Febr. der Königl. Societät vom Hrn. Hofr. \textsc{Gauss} überreichte Vorlesung, überschrieben
\begin{center}
\textit{Theoria combinationis observationum erroribus minimis obnoxiae, pars posterior,}
\end{center}
\noindent steht im unmittelbaren Zusammenhange mit einer frühern, wovon in diesen Blättern [1821 Februar 26] eine Anzeige gegeben ist. Wir bringen darüber nur kurz in Erinnerung, dass ihr Zweck war, die sogenannte Methode der kleinsten Quadrate auf eine neue Art zu begründen, wobei diese Methode nicht näherungsweise, sondern in mathematischer Schärfe, nicht mit der Beschränkung auf den Fall einer sehr grossen Anzahl von Beobachtungen, und nicht abhängig von einem hypothetischen Gesetze für die Wahrscheinlichkeit der Beobachtungsfehler, sondern in vollkommener Allgemeinheit, als die zweckmässigste Combinationsart der Beobachtungen erscheint. Der gegenwärtige zweite Theil der Untersuchung enthält nun eine weitere Ausführung dieser Lehre in einer Reihe von Lehrsätzen und Problemen, die damit in genauester Verbindung stehen. Es würde der Einrich%
